\section{Opérateurs de masse dans HMT--MD : échelle absolue, atlas des parcours et réalisation par secteurs}
\label{sec:tabla-completa-masas-hmt-sm-revfinal}

La masse apparaît en Mécanique Dimensionnelle comme la réalisation de deux branches
mathématiques qui convergent sans se confondre. La première conserve la généalogie de
l'état : une extension survivante détermine un parcours observable ; le parcours écrit
un enregistrement chronologique dodécaphasique ; l'enregistrement produit une signature entière, et la signature
alimente le caractère et la hiérarchie harmonique. La seconde construit la droite
dimensionnelle : le lecteur déterminantal fixe l'échelle d'action, l'horloge interne la
transporte vers l'énergie et la carte cinématique la mène à la masse. L'opérateur d'espèce
agit alors par homothéties positives sur cette droite. La coordonnée expérimentale
n'est introduite qu'après fixation des deux branches.

Le domaine interne complet possède quatre dénombrements distincts :
\begin{equation}
 81=25_{\ell^-}+20_{\nu}+20_d+16_u,\qquad
 56=41\!\times\!1+10\!\times\!2+2\!\times\!3+1\!\times\!4+2\!\times\!5,
 \label{eq:censos-81-56-masas-rectoras}
\end{equation}
\begin{equation}
 13\ \text{multisections},\qquad
 324=81\times4\ \text{routes orientées}.
 \label{eq:censos-13-324-masas-rectoras}
\end{equation}
Les réalisations physiques comparées plus loin constituent une
classification de réalisations ; elles ne sont pas un autre nom pour les 81 cellules, les 56
familles, les 13 multisections ou les 324 histoires orientées. Sur ce domaine,
les deux lecteurs de parcours produisent quatorze profils \(K_D\), neuf profils
\(K_\Omega\), quarante paires conjointes et dix-huit coïncidences. Le secteur nul
demeure distinct du caractère positif de masse.

\subsection{De l'incompatibilité APP à la signature de masse}

La différence primordiale entre les feuillets somme et produit apparaît dans la
décomposition spectrale du bloc central d'APP :
\begin{equation}
 B_c=7I+2R=9P_++5P_-,\qquad
 \operatorname{Spec}(B_c)=\{9,5\},\qquad
 \Delta_{\rm APP}=9-5=4.
 \label{eq:salto-app-cuatro-ley-masas}
\end{equation}
Les entiers qui réapparaissent ensuite conservent des sources distinctes. Le coefficient
local d'échange est \(2\) dans \(B_c=7I+2R\) ; la compression modulaire donne
\(51-45=6\), et son déploiement sur neuf blocs produit \(9\cdot6=54\). Ces
égalités appartiennent à des niveaux typés d'une même généalogie et ne constituent pas
une chaîne de définitions du saut \(4\). Dans la section électronique, \(54\) est
en outre \(D_1(\Gamma_e)\) et la périodicité cinématique d'ordre \(54\) de
l'endomorphisme temporel CT108 : il mesure une composante de l'histoire déjà transportée
par TPK, et non la séparation spectrale primordiale.

Pour chaque parcours enrichi \(\Gamma\), l'enregistrement dodécaphasique détermine
\begin{equation}
 \sigma(\Gamma)
 =(n_A,s_C,k_{\Delta_4},b_{120},b_\zeta)\in\mathbb Z^5,
 \qquad \zeta_{270}=135\Delta_4^2.
 \label{eq:firma-z5-masas-rectoras}
\end{equation}
Le caractère central et son raffinement harmonique sont
\begin{align}
 \log\chi_0(\sigma)
 &=n_A A_{\rm rad}+\frac{s_C}{6}C^*_{\rm rad}
   +k_{\Delta_4}\Delta_4+b_{120}s_{120}+b_\zeta\zeta_{270},
 \label{eq:caracter-central-masas-rectoras}\\
 \log\chi_{\rm full}(\sigma,\eta)
 &=n_A A_{\rm rad}+\frac{s_C}{6}C^*_{\rm rad}
   +k_{\Delta_4}\Delta_4+b_\zeta\zeta_{270}
   +(b_{120}+\eta_{120})s_{120}
   +\sum_{\substack{h\in\mathfrak S\\h\ne120}}\eta_hs_h.
 \label{eq:caracter-completo-masas-rectoras}
\end{align}
La division \(s_C/6\) procède de l'intégralité dodécaphasique. La paire étiquetée
\((b_{120},\eta_{120})\) conserve sa généalogie même lorsque son évaluation est
regroupée en un seul coefficient.

Les moments qui interviennent dans
\eqref{eq:caracter-completo-masas-rectoras} appartiennent à une hiérarchie infinie :
\begin{equation}
 q_\pm=\exp\!\left[-\frac{\pi}{180}(A\pm C^*)\right],
 \qquad s_h=q_-^h-q_+^h,\qquad c_h=q_-^h+q_+^h,
 \label{eq:momentos-torres-masas-rectoras}
\end{equation}
\begin{equation}
 \mathfrak S=\langle90,120\rangle
 =\{90,120\}\cup\{30r:r\ge6\}.
 \label{eq:semigrupo-global-torres-masas-rectoras}
\end{equation}
\(s_n\) comme \(c_n\) satisfont
\begin{equation}
 X_{n+2}=c_1X_{n+1}-(q_-q_+)X_n.
 \label{eq:recurrencia-torres-masas-rectoras}
\end{equation}
Les hauteurs \(90,120,180,210,240,270,360,420\) sont des témoins généalogiques
particulièrement visibles ; le domaine de la loi est le semi-groupe complet. En
particulier, \(360\) conserve deux histoires avant le quotient, quatre pas de
\(90\) et trois de \(120\), même lorsque les deux produisent le même moment scalaire.

\subsection{La section absolue de masse et la décennie \texorpdfstring{\(10^{-34}\)}{10⁻³⁴}}

Le bloc local de Hadamard
\[
 H_4=
 \begin{pmatrix}
  1&1&1&1\\
  1&1&-1&-1\\
  1&-1&1&-1\\
  1&-1&-1&1
 \end{pmatrix},
 \qquad
 \operatorname{SNF}(H_4)=\operatorname{diag}(1,2,2,4),
\]
a un conoyau d'ordre \(16\). Son lecteur déterminantal produit
\begin{equation}
 \Sigma_H=\varphi\alpha_{\rm HMT}^{16},
 \qquad
 10^{-34}<\Sigma_H<10^{-33},
 \qquad
 \left\lfloor\log_{10}\Sigma_H\right\rfloor=-34.
 \label{eq:orden-menos-34-tabla-masas}
\end{equation}
La conversion en \(\mathrm{J\,s}\) constitue une carte ultérieure : la puissance
\(10^{-34}\) naît dans la structure discrète et fixe l'échelle d'une unique droite
d'action. Ses sections antérieure et postérieure au retour sont
\begin{equation}
 \hbar_{\rm pre}=H_5\,10^{-34}\mathcal U_S,\qquad
 \hbar_{\rm ret}=\eta_{\rm ret}\,10^{-34}\mathcal U_S,
 \qquad
 R_{\rm act}=\frac{H_5}{\eta_{\rm ret}}>0.
 \label{eq:accion-absoluta-y-razon-tabla-masas}
\end{equation}
Le dédoublement possède une carte additive,
\(\delta_{2w}=H_5-\eta_{\rm ret}\), et la carte multiplicative adimensionnelle
\(R_{\rm act}\). Une durée interne \(t_0\) étant déclarée, la périodicité de
108 itérations de l'endomorphisme temporel CT108 construit
\begin{align}
 \omega_0&=\frac{2\pi}{108t_0},
 &\varepsilon_{\rm clk}&=\hbar_{\rm ret}\omega_0,
 \label{eq:cuanto-reloj-masa-rectora}\\
 \mathbf m_\star^{\rm pre}
 &=\left(H_5\,10^{-34}\mathcal U_S\right)
   \frac{2\pi}{108t_0}c_{\rm int}^{-2},
 &\mathbf m_\star^{\rm ret}
 &=\left(\eta_{\rm ret}\,10^{-34}\mathcal U_S\right)
   \frac{2\pi}{108t_0}c_{\rm int}^{-2}\in\mathcal L_M.
 \label{eq:seccion-reloj-masa-rectora}
\end{align}
Les deux sections satisfont
\(\mathbf m_\star^{\rm pre}/\mathbf m_\star^{\rm ret}=R_{\rm act}\).
La décennie \(10^{-34}\), l'horloge commune et \(c_{\rm int}^{-2}\) fixent l'échelle
absolue. Dans un rapport de masses, cette section commune se simplifie exactement ; le
raffinement relatif relève du parcours, et non d'une seconde introduction de la
puissance décimale.

\subsection{Lecteurs bidirectionnels de trajectoire et loi opératorielle sur les \(13\) multisections}

Les lecteurs qui agissent sur un parcours \(\Gamma\) sont
\begin{equation}
 K_D(\Gamma)
 =\left(D_{27}+D_{36},\,D_1,\,\frac{D_4-D_3}{2}\right),
 \qquad
 K_\Omega(\Gamma)=(\Omega,54,P_6).
 \label{eq:lectores-kd-komega-masas-rectoras}
\end{equation}
Chaque lecteur induit sur une multisection \(F\) un opérateur diagonal
\(\widehat K_F^r\), \(r\in\{D,\Omega\}\). L'indice \(s\) désigne une section ou un
bloc physique au sein de \(F\), et \(\widehat B_{F,s}^r\) est la restriction
autoadjointe de \(\widehat K_F^r\) à ce bloc. Si
\(\widehat{\mathcal A}^{(0)}_{F,s}\) est la restriction de l'opérateur réunissant la
signature, le caractère complet et la hiérarchie globale des tours, la loi générale prend
la forme
\begin{equation}
 \boxed{
 \widehat{\mathbf M}^{\,\beta,r}_{F,s}
 =\mathbf m_\star^{\rm ret}\otimes
  R_{\rm act}^{\widehat B_{F,s}^r/2}
  \widehat{\mathcal A}^{(0)}_{F,s}
  R_{\rm act}^{\widehat B_{F,s}^r/2},
 \qquad r\in\{D,\Omega\}.}
 \label{eq:ley-operatorial-general-masas-rectoras}
\end{equation}
L'écriture symétrique préserve la positivité sans imposer la commutation de
\(\widehat{\mathcal A}^{(0)}_{F,s}\) et \(\widehat B_{F,s}^r\). Dans la base
simultanément diagonale des certificats en vigueur, elle coïncide avec le produit
\(\mathbf m_\star^{\rm ret}\otimes\widehat{\mathcal A}^{(0)}_F
\exp((\log R_{\rm act})\widehat K_F^r)\).
Cette équation régit les 81 cellules, les 56 familles, les 13 multisections et les
324 parcours orientés. Hors d'une fibre de rang un, le résultat primaire est
l'opérateur complet et son spectre. La réalisation scalaire utilise le morphisme établi
\(\mathcal S_{\rm phys}:(\sigma_{\rm ud},g,\chi_{\rm fina},[\gamma]_{\rm col})
\mapsto\gamma^{\rm surv}_{P,\mathfrak M}\) de
\cref{eq:realizacion-fisica-sabor-ruta-cerrada-rev2} ; cette sélection est
indépendante des masses cibles et précède toute comparaison métrologique.

\subsection{Réduction de rang \(1\) dans la construction de l'électron central}

La section électronique est l'unique point fixe central de l'atlas :
\begin{equation}
 (r,c)=(5,5),\qquad
 \Gamma_e=(5,5,++),\qquad
 \tau_e=\mathtt{+++---+++---}.
 \label{eq:ruta-electronica-central-masas-rectoras}
\end{equation}
Les deux lecteurs y coïncident,
\begin{equation}
 K_D(\Gamma_e)=K_\Omega(\Gamma_e)=(80,54,6),\qquad
 \kappa_e=80-54\alpha_{\rm HMT}+6\Delta_4.
 \label{eq:kappa-electron-rector}
\end{equation}
La réduction de rang un de
\eqref{eq:ley-operatorial-general-masas-rectoras} fournit
\begin{equation}
 \mathbf m_e^\beta
 =\mathbf m_\star^{\rm ret}\,
   \mu_e^{(0)}R_{\rm act}^{\,80-54\alpha_{\rm HMT}+6\Delta_4},
 \label{eq:electron-rank1-ley-masas}
\end{equation}
dont la coordonnée dans la carte comparative est
\begin{equation}
 \boxed{m_e^\beta
 =0.5109989506900008086082571648\ldots\ \mathrm{MeV}/c^2.}
 \label{eq:masa-electron-tabla-completa-revfinal}
\end{equation}
La valeur de base
\(0.5109989224880660725\ldots\ \mathrm{MeV}/c^2\) conserve sa fonction
d'étape intermédiaire ; le résultat bêta de
\eqref{eq:masa-electron-tabla-completa-revfinal} régit la comparaison finale.

\subsection{Classification physique et comparaison quantitative par secteurs}

L'exposé s'organise en quatre strates coordonnées. La première présente le
résultat interne et son codomaine ; la deuxième déclare le morphisme de réalisation
physique ; la troisième montre la meilleure coordonnée numérique déjà disponible et conserve
ses hypothèses et son domaine ; la quatrième introduit la donnée externe comme confrontation ultérieure.
Chaque colonne conserve ainsi son type mathématique : résultat HMT, morphisme de
réalisation, coordonnée calculée et confrontation externe. Le tableau présente
simultanément les évaluations disponibles et les applications de construction qui les
produisent.

L'électron donne sa clôture bêta scalaire. Le photon et les gluons donnent
des sections nulles exactes dans leurs cartes de masse, ainsi que des extensions de jauge non
triviales. Les autres leptons chargés et les six quarks donnent simultanément
leurs parcours nominaux, caractères scalaires, opérateurs diagonaux et spectres sur
leurs fibres ; les neutrinos donnent l'opérateur neutre et son mélange. \(W\), \(Z\) et
\(H\) possèdent des parcours CT--108 individualisés, des empreintes complètes et des caractères
\((n_A,s_C)=(94,-1),(95,-1),(97,9)\). La paire \(K_D/K_\Omega\) appartient au
raffinement des treize multisections de l'atlas et n'est pas une prémisse de la
réalisation bosonique par parcours nominal. Les composés, les secteurs topologiques et les états
virtuels maintiennent respectivement leurs codomaines de liaison, de fusion, de tressage
et d'obstruction.

\Needspace{15\baselineskip}
\subsubsection{Correspondance typée entre objets internes et codomaines physiques}

